\documentclass{article}
\usepackage[T1]{fontenc}

\usepackage{graphicx}
\usepackage{authblk}

% Bibliography (use apacite; don't also load natbib)
\usepackage[natbibapa]{apacite}

% Float control: don't allow figures/tables to appear before they're defined
\usepackage{flafter}
% Optional: stop floats crossing a barrier like a section boundary
% \usepackage[section]{placeins}

% Hyperlinks + smart cross-references (load hyperref late)
\usepackage[hidelinks]{hyperref}
\usepackage[nameinlink]{cleveref}

\providecommand{\keywords}[1]{\textbf{\textit{Keywords:}} #1}

\title{RegCheck: A tool for automating comparisons between study registrations and papers}

\author[* 1,2]{Jamie Cummins}
\author[1]{Beth Clarke}
\author[1]{Ian Hussey}
\author[1]{Malte Elson}

\affil[1]{\small{Institute of Psychology, University of Bern}}
\affil[2]{\small{Bennett Institute of Applied Data Science, University of Oxford}}

\affil[*]{\small{Corresponding author: jamie.cummins@unibe.ch}}

\begin{document}
\maketitle

\begin{abstract}
Across the social and medical sciences, researchers recognize that specifying planned research activities (i.e., ``registration'') prior to the commencement of research has benefits for both the transparency and rigour of science. Despite this, evidence suggests that study registrations frequently go unexamined, minimizing their effectiveness. In a way this is no surprise: manually checking registrations against papers is labour- and time-intensive, requiring careful reading across formats and expertise across domains. The advent of generative language models unlocks new possibilities in facilitating this activity. We present RegCheck, an open-source, modular LLM-assisted tool designed to help researchers, reviewers, and editors from across scientific disciplines compare study registrations with their corresponding papers. More broadly, RegCheck is one implementation of Parity, a general software workflow for automating (meta-)scientific quality-evaluation tasks using the IDEA framework: ingest documents, define evaluation dimensions, extract relevant text, and adjudicate on the basis of that evidence. This framing also makes clear that generative LLMs are used primarily at the adjudication stage, whereas ingestion, dimension definition, and evidence extraction rely on parsers, human input, embeddings, and similarity-based retrieval. Importantly, RegCheck keeps human expertise and judgement in the loop by (i) ensuring that users are the ones who determine which features should be compared, and (ii) presenting the most relevant text associated with each feature to the user, facilitating (rather than replacing) human discrepancy judgements. RegCheck also generates shareable reports with unique RegCheck IDs, enabling them to be easily shared and verified by other users. RegCheck is designed to be adaptable across scientific domains, as well as registration and publication formats. In this paper we provide an overview of the motivation, workflow, and design principles of RegCheck, situate it within the broader Parity/IDEA architecture, and discuss its potential as an extensible infrastructure for reproducible science with an example use case.
\end{abstract}

\keywords{Study registration; large language models; embeddings; research trustworthiness assessment; open source software}

\newpage
\section{Introduction}
The importance of registering study designs, materials, hypotheses, outcomes, and processing/analysis plans for the transparency and rigour of research has seen increasing recognition in the medical and social sciences in the last 20 years. In clinical trials, registration is a legal requirement according to the regulations governing many countries such as the United States \citep{noauthor_nih_nodate} and Europe \citep{noauthor_regulation_2014}. Even when not legally required, registration sees use across the sciences. In general medical research, the creation of registrations (or publication of study protocols) is growing in prevalence \citep{tan_prevalence_2019}; the same trend can also be seen in economics, with pre-analysis plans \citep{ofosu_pre-analysis_2023}, and in psychology, with preregistration (and more recently, Registered Reports, \citeauthor{chambers_past_2021}, 2021; \citeauthor{hardwicke_prevalence_2024}, 2024). Although different fields have different labels, norms, and literatures discussing these approaches, for the sake of consistency here, we refer to all of these using the term ``registration''.

In principle, these registrations improve the transparency of decision making and plans associated with research programs. This transparent record allows researchers to evaluate whether planned analytic choices were followed or not, thus facilitating the trustworthiness of scientific research more broadly \citep{hardwicke_reducing_2023}. Crucially, registrations enable the reader to distinguish which analytic decisions were made before and after observing the data, which can have implications for the veracity and severity of the reported statistical tests \citep{simmons_false-positive_2011}. This provides researchers with critical information that they can use to evaluate the robustness of the results for themselves: in other words, researchers can more accurately assess the risk of bias and calibrate their confidence in the results accordingly \citep{hardwicke_reducing_2023}.

In practice, however, study registration is often more complex. Researchers may need to deviate from their registration for a variety of reasons \cite[see][for a discussion on when and how researchers might consider deviating from their registration]{lakens_when_2024}. Although some deviations could be unproblematic or trivial, others can seriously impact the severity of the test(s) of the predictions \citep{lakens_when_2024}. This is far from a distant concern---several meta-science studies across a variety of fields have found that deviations from registrations are common \citep{bakker_ensuring_2020, claesen_comparing_2021, goldacre_compare_2019, hahn_cross-sectional_2025, heirene_preregistration_2024, li_systematic_2018, ofosu_pre-analysis_2023, poole_systematic_2025, targ_meta-research_group__collaborators_estimating_2023, van_den_akker_selective_2023, van_den_akker_potential_2024}. Even more concerning is that these deviations are often not reported in the final paper \citep{van_den_akker_selective_2023, van_den_akker_potential_2024, targ_meta-research_group__collaborators_estimating_2023}\footnote{Another serious issue is that registrations are often underspecified and do not provide sufficient details to deviate from to begin with \citep{bakker_ensuring_2020, hahn_cross-sectional_2025, ofosu_pre-analysis_2023, poole_systematic_2025, van_den_akker_potential_2024}.}.

We might expect such issues to be caught in the peer review process. Yet, clearly, this is often not the case \citep{van_den_akker_potential_2024}. So how do these deviations make it past peer review? The answer is probably quite simple. Checking registrations against the eventually published paper is a time-intensive and laborious process. Given peer reviewers are already overburdened \citep{adam_peer-review_2025}, the fairly substantial additional task of checking studies against their registrations seems unlikely to be done. Indeed, one study of articles published at \textit{PLOS} suggests that very few editors and reviewers engage with registrations during the peer review process \citep[see also \citeauthor{spitzer_registered_2023}, \citeyear{spitzer_stage_2023}]{syed_data_2025}. In a survey of reviewers for medical journals, only 34.3\% indicated they had examined the information on a trial registry within the past two years \citep{mathieu_use_2013}. This may also serve to explain why deviations are not reported by authors themselves: at least in some cases, it is plausible that authors are unaware of such deviations, given that checking for them is labourious.

While study registration offers a promising solution in principle to the risk of bias that can result from researcher degrees of freedom, in practice, the effectiveness of registration is seriously undermined by the work involved in conducting registration-paper comparisons. If registrations are never examined or evaluated, then we risk assuming that the risk of bias has been reduced when this may not be the case at all.

The recent advent of generative large language models (LLMs) offers the possibility of automating portions of the work involved in registration-paper comparisons, assisting reviewers in a task that they currently seem to avoid (note that we refer to ``registration-paper'' or ``registration-publication'' comparisons throughout this manuscript, but of course such processes can be done with drafts of manuscripts prior to publication or submission to a journal). Indeed, doing this is already accessible to the average researcher: the simplest way to compare registrations and publications is to upload them to a commercial LLM and prompt it to compare them. However, this ad hoc approach raises a host of issues. First, the dimensions of comparison may vary across calls if not explicitly specified: on one call the LLM may compare materials based on sample size, statistical models, and inference criteria, while on another call it may compare the very same materials based on hypothesis tests and primary outcome definitions. While users can specify these dimensions, they may still receive very different outputs depending on their definitions of these dimensions. Additionally, the specific content of prompts will necessarily be unstandardized between researchers, which can introduce a substantial degree of variability in the quality and accuracy of output \citep{cummins_threat_2025,sclar_quantifying_2023}. More generally, the act of verifying the judgements of consistency made by the model may eventually be as laborious as doing the manual comparison in the first place. Most critically, the importance of expert human judgement is displaced by the model's (sometimes erroneous or absent) judgement.

What is needed instead is an approach which ensures consistency in prompt content and comparison logic across calls and users, while also ensuring that expert knowledge is applied by (i) allowing researchers to specify which dimensions should be examined, and (ii) facilitating researchers' assessment of discrepancies, rather than entirely outsourcing this to the model.

More generally, registration-paper comparison is one member of a broader class of (meta-)scientific quality-evaluation tasks in which scientific documents must be ingested, evaluated against explicitly defined dimensions, and judged on the basis of evidence extracted from those documents. We refer to the abstract workflow we use for this class of tasks as Parity. RegCheck is a task-specific implementation of Parity for comparing registrations and papers. Situating RegCheck within Parity helps clarify which parts of the workflow are general, which are specific to registration-paper comparison, and where generative LLMs do and do not enter the process.

To meet these needs, we have developed \href{https://regcheck.app}{RegCheck}: a tool which automatically extracts the most relevant text from registrations and papers as they relate to various user-specified dimensions, easily facilitating comparison between these sources. It provides a pipeline that parses registration and publication texts, lets users specify what dimensions they want to check, and produces structured, auditable reports which are shareable via unique links. By lowering the difficulty barrier for checking registrations without removing expert oversight, RegCheck aims to make consistency checks a routine part of scientific workflows.

Critically, RegCheck represents a single instance of a broader category of emerging meta-scientific tool: namely, one whose purpose is to provide (semi-)automated comparisons between different academic research objects. In this spirit, although this paper is primarily focused on describing RegCheck, it is worth first noting the broader technical infrastructure of which RegCheck represents one branch: namely, an abstract architecture we refer to as Parity. The Parity architecture underlying RegCheck is designed to be adaptable across scientific domains, as well as registration and publication formats. In this paper we provide an overview of the motivation, workflow, and design principles of RegCheck, situate it within the broader Parity architecture, and discuss its potential as an extensible infrastructure for reproducible science with an example use case.

\section{Parity and the IDEA Framework}
Parity is the broader software workflow within which RegCheck is implemented. It is not itself a single end-user tool, but rather an abstract architecture that can be appropriated to implement different (meta-)scientific quality-evaluation tasks in a way that keeps inputs, evidence retrieval, and judgement steps conceptually distinct. The motivation for this abstraction is that many scientific evaluation tasks share the same overall structure even when their objects differ. Registration-paper comparison is one example, but the same logic can in principle apply to tasks such as checking manuscript compliance with reporting guidelines (e.g., ARRIVE, PRISMA), identifying salient information for extraction (e.g., for evidence synthesis), or examining consistency between analysis code and natural language descriptions of analyses within scientific papers.

Parity is organized around a four-step conceptual framework which serves to provide structure in how to think about, break down, and approach (meta-)scientific comparison, extraction, and adjudication question using Parity. We refer to this as the IDEA framework, which consists of the following four core steps:
\begin{enumerate}
\item \textbf{Ingestion.} Collect all documents relevant to the specific task at hand, and parse them in ways that maximise their utility and clarity for that particular task. In practice this may involve interfacing with APIs and handling json responses (e.g., on ClinicalTrials.gov), XML extraction, PDF parsing, document-layout models such as GROBID, reference stripping, section filtering, or other deterministic preprocessing needed to transform scientific materials into task-ready text.
\item \textbf{Dimension Definition.} Specify along which dimensions the scientific materials should be evaluated, compared, or checked. This step is explicitly human-led: users, reviewers, editors, or researchers define the dimensions of interest, whether these are reporting-guideline items, preregistered sample-size commitments, outcomes, or any other target of interest to the user.
\item \textbf{Evidence Extraction.} Represent documents and dimension definitions in a form that supports targeted retrieval, and then extract the most relevant evidence for each dimension. In practice, this is mostly frequently done by the use of semantic embeddings models, which represent semantic information through numeric vectors. Text extraction is then done by splitting scientific documents into chunks, embedding these chunks and dimension definitions specified in the previous step, and querying the similarity of each chunk with the dimension of interest.
\item \textbf{Adjudication.} Make a judgement for the task at hand on the basis of the extracted evidence. This is the stage at which generative LLMs are most naturally used, because the system must synthesize multiple excerpts and render a task-specific judgement, such as whether a registration and paper are consistent on a given dimension.
\end{enumerate}

In addition to providing a scaffolding for researchers to think about how to decompose their meta-scientific workflow into distinct steps, the IDEA framework also helpfully clarifies that different steps in the framework involve different tools or sources of input. Most critically, the IDEA framework disambiguates where generative LLMs are and are not used in these processes. Within the canonical IDEA workflow, Ingestion is primarily a matter of conventional approaches to document parsing, Dimension Definition involves entirely human input, and Evidence Extraction is based on embeddings models. Generative LLMs are by default only invoked at the Adjudication step, where synthesised judgement is actually required. 

This separation of roles makes the workflow easier to inspect, evaluate, and modify, and it helps localize sources of error when they arise. Additionally, for users who are skeptical or reticent to use generative LLMs, this does not render the entirety of Parity workflows inaccessible, since the LLM is only used at the final stage of adjudication, and the preceding stages can be used without it (e.g., one could use the evidence extraction stage to retrieve relevant text excerpts for each dimension, and then make their own human judgement on the basis of that evidence, without using an LLM at all).

RegCheck is therefore best understood as an implementation of Parity which is developed specifically for registration-paper comparison. With this context in mind, the main focus of this paper is to introduce and describe RegCheck itself, in particular its use case, workflow, report structure, and potential value for authors, reviewers, editors, and meta-scientists. Critically, with the scaffolding of Parity and the IDEA framework in mind, this should serve to clarify to readers that RegCheck is not merely ``an LLM that reads and compares two documents,'' but is instead a modular workflow that takes a systematic approach to the task of registration-paper comparison.

\section{The RegCheck Workflow}
RegCheck operationalizes the four IDEA steps through a six-stage backend pipeline. These additional stages reflect specific modifications to the framework that serve to better facilitate the specific task of registration-paper comparison. Stages 1-2 correspond primarily to \textbf{Ingestion}, Stage 4 to \textbf{Dimension Definition}, Stage 3 together with the retrieval component of Stage 5 correspond to \textbf{Evidence Extraction}, and the judgement component of Stage 5 corresponds to \textbf{Adjudication}. Stage 6 is an additional layer focused on reporting, which is not a core part of the IDEA framework but is important to ensure the interpretability and ease of use of the software.

\begin{enumerate}
\item \textbf{Ingestion.} The user specifies a registration (through upload, or by providing a ClinicalTrials.gov identifier) and paper (through upload). Papers are parsed using specialized models of the user's choice (currently GROBID; \citeauthor{grobid_grobid_2025}, 2025, and DPT-2; \citeauthor{huang_document_2025}, 2025, as offered models), while registrations are parsed using standard libraries (PyMuPDF), XML extraction, or API interactions with trial registries, depending on content and source.
\item \textbf{Preparation.} Reference sections are stripped and layout is normalized. If the paper is indicated to consist of multiple studies, then the paper is further processed with irrelevant studies removed, and only the general introduction, the relevant study, and general discussion are retained. When the focal study refers back to details described elsewhere in the paper (e.g., ``Our procedure was the same as the previous experiment''), RegCheck applies additional follow-up extraction logic so that those antecedent details can be surfaced for comparison.
\item \textbf{Embedding and indexing.} All content from both documents is chunked using a token-based, sentence-aware method (200 tokens per chunk, full sentences retained, overlap of 30 tokens between chunks). Embeddings are computed for these text chunks to support retrieval-augmented extraction using OpenAI's \texttt{text-embeddings-3-large} model. These vector representations form the basis of the evidence-extraction process.
\item \textbf{Definition.} Users either select from pre-set dimensions for comparison or define their own by providing a label for the dimension and a definition. Definitions are used for targeted retrieval and for better-specified adjudication prompts.
\item \textbf{Extraction and adjudication.} For each dimension, RegCheck:
\begin{itemize}
    \item Retrieves the most relevant evidence excerpts from each source using dense embedding retrieval. It forms a single query from the dimension label plus its definition (user-provided or a built-in default), embeds this query, and ranks the pre-embedded registration and paper chunks for relevance (based on cosine similarity). The top candidates are then re-ranked with an additional cosine-based scoring step and truncated to the final top-$k$.
    \item Sends the retrieved excerpts to a user-selected LLM, which produces concise summaries of the dimension-relevant registration and paper content (referencing the excerpt IDs).
    \item Produces a deviation judgement from the LLM using the same evidence, encoded as \texttt{yes} (deviation), \texttt{no} (no deviation), or \texttt{missing} (insufficient evidence), plus a brief explanation in \texttt{deviation\_information}.
\end{itemize}
\item \textbf{Reporting.} RegCheck renders a structured results page (and CSV export) keyed by a task ID. For each dimension it displays the retrieved evidence excerpts from the registration and paper (with excerpt IDs and relevance scores), the LLM summaries, and the deviation judgement plus supporting explanation. \Cref{fig:regcheck-report} provides an illustration of the current appearance of the RegCheck report.
\end{enumerate}

\begin{figure}[!htbp]
    \centering
    \includegraphics[width=\linewidth]{RegCheck-report-demo.png}
    \caption{A screenshot of an example RegCheck report. Each row represents a comparison between the registration and paper along a specific, pre-specified dimension (noted in the ``Dimension'' column). Under the ``Preregistration'' and ``Paper'' columns, users can toggle between displaying direct quotes from the respective document, or LLM-generated summaries of the content of these quotes. The color-coding of the row indicates whether a deviation between the materials was found (red), no deviation was found (blue), or that there was insufficient information in one or both sources to render a deviation judgement (yellow). The ``Deviation Information'' column also provides commentary on the similarities and differences between the two sets of materials.}
    \label{fig:regcheck-report}
\end{figure}

\section{Design Principles and Philosophy}
\begin{itemize}
\item \textbf{Pragmatism.} RegCheck is a software tool which provides a standardized framework for checking registrations against papers, with the goal of increasing the frequency of registration-paper comparisons (which is not currently done regularly). All steps of its research and development are conducted with the aim of evaluating and improving its extraction accuracy, informativeness, and ease of use.
\item \textbf{Human-in-the-loop.} RegCheck facilitates human experts; it does not replace them. The system highlights candidate discrepancies, and also provides indicative LLM-based judgements, but it is not aimed at outsourcing human expertise.
\item \textbf{Non-prescriptive.} It is not the role of RegCheck to make determinations about which dimensions matter when comparing a registration to a paper: this can vary by user, use case, and field/discipline norms. There is currently no cross-discipline consensus, that we know of, on the selection of dimensions which do or do not matter for study registration. RegCheck provides default dimensions based on frequently examined features (e.g., sample size, primary outcomes), but also provides free-text fields for users to add their own for this reason. The current default dimensions are most relevant to psychology, but future iterations of the software will offer recommended default dimensions which vary based on the discipline of the to-be-compared paper.
\item \textbf{Discipline-agnostic and discipline-aware.} Although it has origins in psychological research, RegCheck is an abstract workflow aimed at being useful for researchers across a range of scientific disciplines. At the same time, we recognize that RegCheck's accuracy and usefulness may vary across fields and disciplines; for this reason, we are investigating variation in its accuracy in a range of different areas.
\item \textbf{Modular separation of roles.} Because RegCheck is built as a Parity implementation, it explicitly separates document ingestion, human definition of dimensions, similarity-based evidence extraction, and model-based adjudication. This makes the system easier to audit, evaluate, and extend, and it clarifies that errors in one part of the workflow need not be attributed wholesale to ``the LLM''.
\end{itemize}

\section{Privacy and Confidentiality}
RegCheck does not request, log, or store any text uploaded by users, with the exception of extracted quotes from the manuscript during the creation of the RegCheck report (which are stored to facilitate the creation and maintenance of the persistent and shareable RegCheck report pages). Additionally, no user information is logged by RegCheck.

We present users with optional, anonymous surveys upon running a RegCheck report, which we use to improve user experience and provide us with information about the purposes for which the platform is used (e.g., for peer review, self-review, etc.). However, these data are not identifiably linked to the content uploaded by those users.

RegCheck provides users with the option to decide which generative LLM is used in the creation of the RegCheck report (by default, ChatGPT-5 is selected). Some of these models (e.g., DeepSeek, Llama) are open-weights models which are hosted on third-party inference API providers (e.g., the Groq service). Others (e.g., ChatGPT) are closed-source models which are hosted by those commercially invested in that model (e.g., OpenAI serves the ChatGPT model). OpenAI states in their documentation that text served to its models via its API service is not retained nor used for training. However, we of course have no control over the enforcement of, and adherence to, these privacy policies, nor do we have the capacity to know whether such policies may change in the future. We therefore provide the open-weights models as alternative options, which are hosted by third-party providers who are not commercially invested in the development of the models they host. In the future, we have plans to migrate our services to an internal university server, which will also allow us to deploy university-hosted open-weights models in order to provide more definitive guarantees around privacy, confidentiality, and future reproducibility.

Usage of RegCheck should be constrained by the same standards of confidentiality and privacy around academic materials that apply in any other professional context. In other words: just as one should not share confidential materials with unauthorized colleagues, one should not upload such materials to RegCheck. More generally, RegCheck should categorically not be used for content that users do not have permission to share (e.g., materials which users have acquired during the process of peer review which should not be shared beyond reviewers themselves).

\section{Is RegCheck Available To Use?}
Yes, RegCheck is available to use. The cost of running the RegCheck processes varies by the size of registrations and papers and the selected model, but costs on average about 40 cents per report using ChatGPT-5 with medium reasoning effort. RegCheck can be used either through our point-and-click web app (described below) or run locally (at one's own cost) using the open-source code, available on \href{https://github.com/JamieCummins/regcheck}{Github}. When using the public app, the cost of usage is currently covered by our research group; in other words, using RegCheck is entirely free for users.

RegCheck's point-and-click production application is currently available at \href{https://regcheck.app}{https://regcheck.app}. Given that systematic, formal evaluation is currently under way, we currently do not provide any guarantees about the fidelity of RegCheck's output. Indeed, we strongly advise that users manually verify all output for accuracy. With this said, our impressions of RegCheck, based on extensive but unsystematic testing, are that it is generally very accurate in its extractions and is quite adept at catching genuine deviations between registrations and papers where these deviations exist. RegCheck frequently catches deviations that members of our research team missed, including subtle differences in registered vs. reported numeric inclusion criteria thresholds, changes in planned analyses, and instances of outcome-switching in clinical trials.

\section{When and How Should I Use RegCheck?}
The most obvious use case for RegCheck is for editors and reviewers who are assessing a manuscript where one or more studies were registered (where they have permission to share and upload those materials; see Privacy and Confidentiality). Similarly, research synthesists might use it when assessing the trustworthiness of studies included in a systematic review or meta-analysis. Of course, RegCheck is also useful for authors who want to check their own manuscripts for any deviations they should report \cite[e.g., RegCheck could help to spot registration deviations that might be missing in a registration deviation table, like the one suggested by][]{willroth_best_2024}. Regardless of whether the user is using RegCheck to assess their own paper or someone else's, the process is the same.

RegCheck is designed to be intuitive and easy to use. First, users navigate to \href{https://regcheck.app}{regcheck.app} (currently in beta). The user can then open the ``Tools'' menu of the navigation bar and navigate to ``Registration-Paper Comparisons''. The user may then choose to either use our recommended default settings, or decide for themselves which settings they wish to use. If users opt to choose settings for themselves, they are provided with several parameters to decide on:
\begin{enumerate}
  \item Parser: Users may decide which engine is used for parsing the uploaded paper. Options currently provided are GROBID \citep{grobid_grobid_2025}, which is best suited for papers already typeset to journal standards, and DPT-2 \citep{huang_document_2025}, which is best suited for non-typeset manuscripts (e.g., drafts of manuscripts). We recommend GROBID as the default.
  \item Language model: Users may select the large language model which is used for rendering judgements with regard to deviations in the manuscript. Currently we recommend ChatGPT-5 on medium reasoning effort. There is currently no option for users to select the embeddings model used for retrieval: this is, by default, \texttt{text-embeddings-3-large} from OpenAI \citep{openai_openai_nodate}.
  \item Comparison dimensions: Users may determine which dimensions the registration and paper should be compared along. The default dimensions we suggest are based on our understanding of those factors which many researchers who conduct such comparisons are interested in (e.g., outcome-switching, sample-size specification, hypothesis specification).
  \item Append outputs: Users may determine whether the evaluation of a specific dimension should include the content of previously evaluated dimensions or not. One of the crucial advantages of RegCheck is that it allows for different comparison dimensions to be evaluated separately in separate calls to the LLM. This means that the quality of output for each dimension can be evaluated on an individual basis, since the calls made are independent\footnote{By contrast, if multiple dimensions were evaluated based on a single call to the model, this would mean that individual dimension results were necessarily affected by the output of other dimensions, due to the nature of next-token prediction in language models; by extension, there would be an inherent dependency among outputs for each dimension, which is suboptimal for evaluation.}. It is also possible that one may wish for the software to be context-aware of the content of response(s) to previous dimensions when evaluating another dimension. For example, for clinical trials, knowing what output was previously given for primary outcome definitions might help the software give more accurate output related to secondary outcomes, as well as reduce the likelihood that it confuses a primary outcome for a secondary outcome (i.e., since it is now provided with the context that certain outcomes have already previously been defined as primary). Our recommended default is to allow for this contextual awareness.
\end{enumerate}
After providing this information, users can provide both the registration (either by uploading a Microsoft Word \texttt{docx} or a PDF, or by providing a ClinicalTrials.gov identifier) and the paper (\texttt{docx} or PDF). PDF documents will usually work best for this, as they can be parsed more easily. Users can then simply click ``Compare''.

RegCheck then generates a report. Depending on the model selected and number of dimensions specified, this usually takes anywhere from 1 to 10 minutes. The RegCheck report will populate iteratively as each dimension is evaluated and output is provided. A status message is provided to users specifying how many dimensions have been evaluated thus far, and how many dimensions remain. The final, completed report extracts direct quotes from both the registration and paper with regard to the dimension compared, as well as a judgement of whether the registration and paper match on this dimension (alternatively, RegCheck flags if the relevant information is missing or not reported), and details related to this decision. This output serves an orienting function to users: where before, one would need to scan each line of a registration and paper to compare dimensions, RegCheck provides direct quotes from each document which can be used as search terms to easily identify the relevant content in context in the manuscript (e.g., by using Ctrl+F in the manuscript PDF and searching for the extracted quote text).

RegCheck's provided judgement of deviations, and deviation information, can additionally be useful for spotting deviations that may have been overlooked. However, RegCheck is not intended to replace human judgment. Users must apply their own discretion when evaluating registration-paper consistency, and they should be aware that RegCheck can make both false positive errors (cases where RegCheck flags that there is a deviation when there is not a deviation; also keep in mind that very minor deviations may be flagged) and false negative errors (cases where RegCheck flags that there is no deviation when there is a deviation). We are currently in the process of assessing when, and how frequently, RegCheck makes these errors, as well as validating RegCheck more broadly.

\section{How Accurate Is RegCheck?}
As mentioned, the anecdotal testing of RegCheck among our research team and close colleagues gives us the impression that RegCheck is adept at flagging (in)consistencies between registrations and papers. However, anecdotal testing is not a substitute for formal evaluation. Our research team is thus currently working on several different projects aimed at quantifying and evaluating the performance of RegCheck for comparisons of registrations and papers in psychology, medicine, and economics. Evaluating RegCheck comes with challenges. Importantly, there is little consensus on how one can assess such performance in humans, let alone automated tools. This is largely due to two issues. First, there is no agreed-upon framework on how to evaluate registration-paper consistency: different researchers disagree about what should be specified in a registration, what dimensions along which a registration should be compared to a paper, and what ultimately counts as a deviation or not. Second, some aspects of comparison are inherently ambiguous: whereas it can be relatively straightforward to get agreement on a prespecified sample size, other features of comparison (e.g., prespecified hypotheses or sufficiency of prespecified outcome descriptions) can elicit more variation in the evaluations of humans.

Put simply: in most cases here, establishing a ground truth among human comparators is incredibly difficult, and we lack systematic frameworks for making such comparisons in the first place. In evaluating RegCheck, we acknowledge this fact, and use it to inform our evaluation approach. Specifically, rather than trying to evaluate RegCheck as a software aimed at approximating a known ground-truth value, we instead treat and evaluate it as an accompanying reviewer. In practice, this means one of our core evaluation metrics of interest is the degree of agreement with other human comparators, and the establishment of relative non-inferiority to observed human-human agreement. In other words: for us to consider RegCheck accurate, it should, in general, show agreement with human raters to at least the same extent as humans show agreement with one another. More concretely, our goals in evaluating RegCheck are in terms of interrater agreement on (i) extracted text from registrations/papers (e.g., planned research questions, planned statistical tests), and (ii) judgements of consistency between the registration and paper. To do this for a given research area, we specify and define a set of dimensions along which registrations and papers should be compared, solicit multiple human comparators to provide these comparisons, and examine agreement among these human evaluations and those of RegCheck. In cases of disagreement, we will also examine whether RegCheck is hallucinating or making mistakes, or capturing legitimate details which humans tend to miss.

Although our primary investigations focus on the aspect of RegCheck-human consistency, we additionally plan to examine RegCheck's accuracy through artificially generating registration and paper texts using a generative language model which contain specific injected inconsistencies. Such an approach is more amenable to conventional means of evaluation, given that we are in control of the ground-truth consistency in these cases. This approach also provides us with opportunities to conduct evaluation at greater scales: whereas soliciting human comparators to review registration-paper consistency takes a substantial deal of time, evaluating RegCheck's performance on artificially generated examples can be done much more quickly and efficiently. Of course, the drawback to this is that the artificial texts may be relatively generic, and written in ``LLM-speak'' which will likely not be representative of the diversity of writing styles present in the true academic literature. Consequently, these diverging approaches (interrater agreement with human comparators, artificially generated texts with known (in)consistencies) provide us with multiple perspectives on the accuracy of RegCheck: both its conventional accuracy in an artificial context, as well as its accuracy in more realistic settings.

Of course, in principle RegCheck could still serve a pragmatic function even if it exhibits lower interrater agreement than that observed for human-human agreement. After all, the baseline rate of conducting registration-paper comparisons in typical human reviewing is low, and this is largely due to the fact that doing such comparisons is time-consuming and laborious. Checking the accuracy of an existing RegCheck report, by contrast, is much less time-consuming. In this sense, RegCheck may serve an orienting function to authors and reviewers. Even if its extractions are at times inaccurate, it may functionally increase the likelihood of these checks being done in the first place, where currently the baseline rate of this behavior is very low. However, this is obviously not optimal: our goal with RegCheck is to provide refined software whose text extractions and consistency judgements reach the same standard of agreement with humans as those of another human.

\section{Future Agenda}
RegCheck is best seen as infrastructure for reproducible science. By making registration-paper comparison faster, more consistent, and more transparent, we hope to increase the likelihood that these checks become routine in authors' own workflows as well as in editorial workflows, peer review, and post-publication auditing. Its modularity allows adaptation across disciplines and integration into larger ecosystems for research integrity.
The research arm of RegCheck's development focuses on concretely addressing:
\begin{enumerate}
\item The accuracy of RegCheck's text extraction, and whether this accuracy varies across domains, disciplines, and model hyperparameters;
\item Whether providing RegCheck to authors can improve the quality of registration-paper comparisons upon submission of papers to journals;
\item Whether RegCheck improves the efficiency and frequency of registration-paper comparisons.
\end{enumerate}
We welcome engagement and input from researchers, journals, and funders interested in using RegCheck in diverse research areas. We are currently working to develop extensions of RegCheck, including the capacity to (i) assess multiple papers at once for use in meta-science or meta-analysis research, and (ii) assess consistencies and deviations between study analysis code and descriptions of planned analyses.

\bibliographystyle{apacite}
\bibliography{references}

\end{document}
